\clearpage
\section{Gaussian channels}\label{sec:channels}

\noindent\motiv{\reve{The automata of the previous chapters are reversible: each step is an
automorphism of the CCR algebra and can be undone by its inverse. A lattice of oscillators coupled
to an environment loses this property. Its dynamics is still a physical operation on the
observables, that is, a unital completely positive map, but the map need not have a completely
positive inverse, and in general it adds noise to the states. When the environment is Gaussian and
the coupling quadratic, such a map is expected to send Weyl operators to multiples of Weyl
operators, as the dilations later in the chapter indicate: a Weyl operator $W(u)$ goes to a
multiple of $W(Lu)$, where the linear map $L$ need no longer preserve $\sigma$. The multiple consists of a phase $e^{\ii\mu(u)}$, which only displaces the states, and a
damping factor $e^{-\nu(u,u)/4}$ with a quadratic form $\nu$, the noise, which compensates the loss
of symplecticity. This chapter develops these Gaussian channels over an arbitrary real vector space
with an alternating form. The lattice enters only at the end, where we extend the definition of a quantum
cellular automaton to completely positive dynamics and introduce two notions of locality; Gaussian automata on the chain \reve{are treated in Chapter~\ref{sec:chain}}.}}

\medskip
\noindent\motiv{\reve{Throughout, $V$ is a real vector space with an alternating bilinear form
$\sigma$, no longer only the lattice phase space, and $\CCR(V,\sigma)$ is the C*-algebra generated by
the Weyl operators $W(u)$ on $\ell^2(V)$, constructed exactly as in Chapter~\ref{sec:ccrautomata}, with
$W(u)W(w)=e^{-\ii\sigma(u,w)/2}W(u+w)$ and $W(u)^*=W(-u)$. The letters $V$, $\sigma$ and $L$ thus
have a more general meaning in this chapter than in the table of notation, the lattice data being
one instance. \reve{We define states, characteristic functions and Gaussian states on $\CCR(V,\sigma)$
as in Chapter~\ref{sec:gaussianstates}, whose definitions, and the propositions cited from
it below, use no property of the lattice.} For bilinear
forms $\gamma$ and $\beta$ on $V$ we write $\gamma+\ii\beta\ge0$ if the matrix
$\bigl[\gamma(u_j,u_k)+\ii\beta(u_j,u_k)\bigr]_{j,k}$ is positive semidefinite for every finite family
$u_1,\dots,u_m$ in $V$; for $\beta=\sigma$ this is the state condition of
Chapter~\ref{sec:gaussianstates}.}}

\subsection{Completely positive maps}\label{sec:ch-cp}

\noindent\motiv{\reve{A linear map $T$ between C*-algebras is \emph{positive} if it maps positive
elements to positive elements, and \emph{completely positive} if, for every $k$, the map applied
entrywise to $k\times k$ matrices over the first algebra is positive. Positivity refers to the order
of a C*-algebra, in which the positive elements are the self-adjoint elements with spectrum in
$[0,\infty)$; by a basic theorem of the subject these are exactly the elements of the form $x^*x$,
and finite sums of such are again of this form~\cite{bratteli-robinson} \vetsource{(pinpoint to be
checked)}. The CCR algebra is a closed subalgebra of the bounded operators on a Hilbert space,
and we first check that the order it inherits from there is this intrinsic one. Applied also to the
algebra of $k\times k$ matrices over $\CCR(V,\sigma)$, a closed subalgebra of the bounded operators on
the $k$-fold direct sum of the Hilbert space, this makes complete positivity of maps of the CCR
algebra unambiguous.}}

\begin{proposition}[Positivity in closed subalgebras]\label{prop:ch-starordered}
\reve{Let $A$ be a C*-algebra and $B\subseteq A$ a norm-closed $*$-subalgebra containing the unit of
$A$. An element $x\in B$ is positive in $A$ if and only if $x=s^*s$ for some $s\in B$. In particular
an element of $\CCR(V,\sigma)$ is positive in the C*-algebra $\CCR(V,\sigma)$ if and only if it is a
positive operator.}
\end{proposition}

\begin{proof}
\reve{If $x=s^*s$, then $x$ is positive in $A$ by the theorem just quoted. Conversely, let $x\in B$ be
positive in $A$, so that $x$ is self-adjoint with spectrum in a bounded interval $[0,c]$. By the
Weierstrass theorem there are real polynomials $p_j$ converging to the square root uniformly on
$[0,c]$, and the continuous functional calculus of $A$ gives $p_j(x)\to s=\sqrt x$ in norm. Each
$p_j(x)$ lies in $B$, which contains the unit and is closed, so $s\in B$; it is self-adjoint, and
$s^*s=x$. \reve{For the last statement take for $A$ the bounded operators on $\ell^2(V)$ and
$B=\CCR(V,\sigma)$: an element is positive in $A$ exactly when it is a positive operator, and by the
theorem quoted above, applied in $B$, it is positive in $B$ exactly when it has the form $s^*s$ with
$s\in B$, which by the first statement is the same condition.}}
\end{proof}

\begin{proposition}[Composition of completely positive maps]\label{prop:ch-cpcomp}
\reve{The identity map is completely positive, and the composite of two completely positive maps is
completely positive.}
\end{proposition}

\begin{proof}
\reve{Applying $g\circ f$ entrywise to a matrix is the same as applying $f$ entrywise and then $g$,
and each step preserves positivity.}
\end{proof}

\noindent\motiv{\reve{A positive map between C*-algebras is automatically bounded, and every
$*$-homomorphism is completely positive~\cite{paulsen}. We use both facts below: the first to extend identities
from Weyl operators to the whole algebra, the second to regard automorphisms as channels.}}

\subsection{Gaussian channels and the product formula}

\noindent\motiv{\reve{A Gaussian channel is specified by the data described at the beginning of the
chapter, a linear map, a noise form and a displacement, together with a completely positive map that
realizes them on Weyl operators. The data are subject to a positivity condition, in which the noise
compensates the failure of the linear map to preserve $\sigma$.}}

\begin{definition}[Gaussian channel]\label{def:ch-channel}
\reve{A \emph{Gaussian channel} on $\CCR(V,\sigma)$ consists of a linear map $L\colon V\to V$, a
symmetric bilinear form $\nu$ on $V$ (the \emph{noise}), a linear form $\mu$ on $V$ (the
\emph{displacement}) and a completely positive map $T$ of $\CCR(V,\sigma)$ into itself, such that
\begin{equation}\label{eq:ch-channel}
T\bigl(W(u)\bigr)=e^{\ii\mu(u)-\nu(u,u)/4}\,W(Lu)\qquad(u\in V)
\end{equation}
and
\begin{equation}\label{eq:ch-cond}
\nu+\ii\bigl(\sigma-\sigma(L\cdot,L\cdot)\bigr)\ge0 .
\end{equation}
The \emph{dissipation form} of the channel is
\[\Delta_T(u,w)=\nu(u,w)+\ii\bigl(\sigma(u,w)-\sigma(Lu,Lw)\bigr).\]}
\end{definition}

\noindent\motiv{\reve{Condition \eqref{eq:ch-cond} says that the dissipation form is positive
semidefinite. The literature calls the maps $T$ of this kind quasi-free completely positive maps. For
finitely many modes it is known that \eqref{eq:ch-cond} is necessary and sufficient for a
completely positive map with \eqref{eq:ch-channel} to exist~\cite{dvv} \vetsource{(scope of the
theorem there to be checked)}; these notes prove neither direction, and both the map and the
condition are part of the definition. Since $T$ is linear and bounded, it is determined by
\eqref{eq:ch-channel}, as the following lemma shows.}}

\begin{lemma}[Unit and uniqueness]\label{lem:ch-unit}
\reve{A Gaussian channel satisfies $T(1)=1$. Two bounded linear maps from $\CCR(V,\sigma)$ into a
normed space that agree on all Weyl operators are equal.}
\end{lemma}

\begin{proof}
\reve{Since $W(0)=1$, \eqref{eq:ch-channel} at $u=0$ gives $T(1)=1$. The linear span of the Weyl
operators is a $*$-algebra, because a product of Weyl operators is a multiple of a Weyl operator, and
it is dense in $\CCR(V,\sigma)$ by definition of the latter; since two bounded linear maps that agree
on a set agree on the closure of its linear span, maps agreeing on all Weyl operators are equal.}
\end{proof}

\noindent\motiv{\reve{A $*$-automorphism maps products to products. A Gaussian channel does so
\reve{only up to a scalar for products of two Weyl operators}, and the scalar is governed by the
dissipation form.}}

\begin{proposition}[Product formula]\label{prop:ch-product}
\reve{For all $u,w\in V$,
\[T\bigl(W(u)W(w)\bigr)=e^{-\Delta_T(u,w)/2}\;T\bigl(W(u)\bigr)\,T\bigl(W(w)\bigr).\]}
\end{proposition}

\begin{proof}
\reve{By the Weyl relation and \eqref{eq:ch-channel}, the left-hand side equals
\[e^{-\ii\sigma(u,w)/2}\,e^{\ii\mu(u+w)-\nu(u+w,u+w)/4}\,W(Lu+Lw),\]
while the product on the right equals
\[e^{\ii\mu(u)+\ii\mu(w)-\nu(u,u)/4-\nu(w,w)/4}\,e^{-\ii\sigma(Lu,Lw)/2}\,W(Lu+Lw).\]
Since $\mu$ is linear and $\nu$ symmetric, $\nu(u+w,u+w)=\nu(u,u)+2\nu(u,w)+\nu(w,w)$, and the ratio
of the two scalars is
\[\exp\Bigl(-\tfrac12\nu(u,w)-\tfrac\ii2\sigma(u,w)+\tfrac\ii2\sigma(Lu,Lw)\Bigr)=e^{-\Delta_T(u,w)/2}.\]}
\end{proof}

\begin{example}[The attenuator]\label{ex:ch-attenuator}
\reve{Let $V=\R^2$ be one mode, $\sigma(u,w)=u^{\mathsf T}\Omega w$, and fix $0<\zeta<1$. The data $L=\zeta\cdot1$,
$\nu(u,w)=(1-\zeta^2)\,u^{\mathsf T}w$ and $\mu=0$ satisfy \eqref{eq:ch-cond}, because the dissipation form is
$(1-\zeta^2)\bigl(u^{\mathsf T}w+\ii u^{\mathsf T}\Omega w\bigr)$, a positive multiple of the vacuum's form
$1+\ii\Omega$, whose eigenvalues are $0$ and $2$. This is the attenuator of quantum optics, which
transmits the fraction $\zeta^2$ of the intensity and fills the loss with vacuum noise; the
beam splitter of Example~\ref{ex:ch-beamsplitter} indicates where a completely positive map realizing it
comes from. We use it as the running example of the chapter. Statements about the attenuator below
concern the Gaussian channel with these data, which is unique by Lemma~\ref{lem:ch-unit} and whose
existence is not proved in these notes.}
\end{example}

\subsection{Special cases and reversible channels}

\noindent\motiv{\reve{Condition \eqref{eq:ch-cond} simplifies in three cases. Without noise it forces
$L$ to preserve $\sigma$; for $L$ preserving $\sigma$ it only asks the noise to be non-negative; and
at the opposite extreme, $L=0$, it is the state condition. In the last case every $T(W(u))$ is a
multiple of the unit, so that the functionals $\omega\circ T$ agree on Weyl operators for all states
$\omega$, and hence, by Lemma~\ref{lem:ch-unit}, the channel replaces every state by one fixed state
(a state by the argument in the proof of Proposition~\ref{prop:ch-states}), with characteristic
function $e^{\ii\mu(u)-\nu(u,u)/4}$.}}

\begin{proposition}[Special cases]\label{prop:ch-special}
\reve{Let $L$ be linear and $\nu$ symmetric.
\begin{enumerate}[label=(\alph*)]
\item If $0+\ii(\sigma-\sigma(L\cdot,L\cdot))\ge0$, then $\sigma(Lu,Lw)=\sigma(u,w)$ for all $u,w$.
\item If $L$ preserves $\sigma$, then \eqref{eq:ch-cond} holds if and only if $\nu(u,u)\ge0$ for all $u$.
\item If $L=0$, then \eqref{eq:ch-cond} is the state condition $\nu+\ii\sigma\ge0$.
\end{enumerate}}
\end{proposition}

\begin{proof}
\reve{\begin{enumerate}[label=(\alph*)]
\item Put $a=\sigma(u,w)-\sigma(Lu,Lw)$. For the pair $u,w$ the matrix of the hypothesis is
$\begin{psmallmatrix}0&\ii a\\-\ii a&0\end{psmallmatrix}$, whose determinant $-a^2$ must be
non-negative; hence $a=0$.
\item The matrices of \eqref{eq:ch-cond} are the Gram matrices $[\nu(u_j,u_k)]$ of the real symmetric
form $\nu$, which are positive semidefinite for all families exactly when $\nu(u,u)\ge0$ for all $u$.
\item For $L=0$ the form $\sigma(L\cdot,L\cdot)$ vanishes, so \eqref{eq:ch-cond} reads $\nu+\ii\sigma\ge0$.
\end{enumerate}}
\end{proof}

\begin{proposition}[Reversible channels]\label{prop:ch-reversible}
\reve{Let $L$ be a linear bijection of $V$ preserving $\sigma$. With $\nu=0$ and $\mu=0$, the
automorphism $\alpha_L$ with $\alpha_L(W(u))=W(Lu)$ (Lemma~\ref{lem:relabel}, whose proof uses only that
$L$ is a bijection preserving $\sigma$) is a Gaussian channel.
Conversely, a Gaussian channel with $\nu=0$, $\mu=0$ and bijective $L$ is $\alpha_L$.}
\end{proposition}

\begin{proof}
\reve{As a $*$-automorphism, $\alpha_L$ is completely positive, \eqref{eq:ch-channel} holds, and the
matrices of \eqref{eq:ch-cond} vanish. Conversely, $L$ preserves $\sigma$ by
Proposition~\ref{prop:ch-special}(a), so $\alpha_L$ is defined; the channel and $\alpha_L$ are bounded
and agree on Weyl operators, hence coincide by Lemma~\ref{lem:ch-unit}.}
\end{proof}

\subsection{Concatenation and the action on states}

\noindent\motiv{\reve{Applying two channels in succession is again a channel. In the Heisenberg
picture the first channel to act on a state is the last to act on an observable, so that $T_2\circ
T_1$ applies $T_1$ to the observable first, while on states $T_2$ acts first,
$\omega\mapsto\omega\circ T_2\mapsto\omega\circ T_2\circ T_1$.}}

\begin{proposition}[Concatenation]\label{prop:ch-comp}
\reve{If $T_1$ and $T_2$ are Gaussian channels, then $T_2\circ T_1$ is a Gaussian channel with linear
map $L_2L_1$, noise $\nu_1+\nu_2(L_1\cdot,L_1\cdot)$ and displacement $\mu_1+\mu_2\circ L_1$. Its
dissipation form is $\Delta_{T_1}(u,w)+\Delta_{T_2}(L_1u,L_1w)$.}
\end{proposition}

\begin{proof}
\reve{The composite is completely positive by Proposition~\ref{prop:ch-cpcomp}, and applying
\eqref{eq:ch-channel} twice gives
\[T_2\bigl(T_1(W(u))\bigr)=e^{\ii\mu_1(u)-\nu_1(u,u)/4}\,e^{\ii\mu_2(L_1u)-\nu_2(L_1u,L_1u)/4}\,
W(L_2L_1u).\]
The formula for the dissipation form follows from the decomposition
\[\sigma(u,w)-\sigma(L_2L_1u,L_2L_1w)=\bigl(\sigma(u,w)-\sigma(L_1u,L_1w)\bigr)
+\bigl(\sigma(L_1u,L_1w)-\sigma(L_2L_1u,L_2L_1w)\bigr).\]
Consequently
the matrix of \eqref{eq:ch-cond} for a family $(u_j)$ is the sum of the matrix of the first channel
for $(u_j)$ and that of the second for $(L_1u_j)$, and a sum of positive semidefinite matrices is
positive semidefinite.}
\end{proof}

\noindent\motiv{\reve{In the Schr\"odinger picture a channel acts on states, sending $\omega$ to
$\omega\circ T$. The following proposition shows that Gaussian states go to Gaussian states, with
the covariance transformed by $L$ and increased by the noise.}}

\begin{proposition}[Action on Gaussian states]\label{prop:ch-states}
\reve{Let $\omega$ be a Gaussian state with covariance $\gamma$ and displacement $d$. Then $\omega\circ T$
is a Gaussian state with covariance $\gamma(L\cdot,L\cdot)+\nu$ and displacement $d\circ L+\mu$.}
\end{proposition}

\begin{proof}
\reve{The functional $\omega\circ T$ is bounded, and it takes the value $1$ at the unit because
$T(1)=1$. For every $b$ in the CCR algebra the element $T(b^*b)$ is positive, since $T$ is, and
therefore of the form $s^*s$, so that $\omega(T(b^*b))=\omega(s^*s)\ge0$; thus $\omega\circ T$ is a
state. Its characteristic function is
\[\omega\bigl(T(W(u))\bigr)=e^{\ii\mu(u)-\nu(u,u)/4}\,e^{\ii d(Lu)-\gamma(Lu,Lu)/4},\]
which is the Gaussian characteristic function with the symmetric form $\gamma(L\cdot,L\cdot)+\nu$ and the linear form
$d\circ L+\mu$.}
\end{proof}

\noindent\motiv{\reve{In the running example, the attenuator carries \reve{a Gaussian state with the
vacuum covariance}, $\gamma(u,w)=u^{\mathsf T}w$ and $d=0$, to the state with covariance
\[\zeta^2u^{\mathsf T}w+(1-\zeta^2)\,u^{\mathsf T}w=u^{\mathsf T}w,\]
so that it leaves the vacuum \reve{covariance} invariant, as it should, since it only mixes the mode
with vacuum noise. A \reve{Gaussian state with the thermal covariance} $\gamma(u,w)=c\,u^{\mathsf T}w$,
$c>1$, goes to the state with covariance $(\zeta^2c+1-\zeta^2)\,u^{\mathsf T}w$, so that the excess $c-1$ over the vacuum is multiplied
by $\zeta^2$.}}

\subsection{Finitely many modes}

\noindent\motiv{\reve{For $n$ modes, $V=\R^{2n}$ with $\sigma(u,w)=u^{\mathsf T}\Omega w$, Gaussian states
are described by covariance matrices, and the covariance matrix $\gamma_B$ of a covariance form $\gamma$
is $\Omega^{\mathsf T}\gamma\Omega$ (Chapter~\ref{sec:gaussianstates}); it is admissible if $\gamma_B^{\mathsf
T}=\gamma_B$ and $\gamma_B-\ii\Omega\ge0$. In this language the data $(L,\nu)$ of a Gaussian channel
become a pair of matrices, while the displacement is carried separately.}}

\begin{definition}[Channel matrices]\label{def:ch-matrices}
\reve{Real $2n\times2n$ matrices $(X,Y)$ are \emph{channel matrices} if $Y^{\mathsf T}=Y$ and
$Y-\ii\Omega+\ii X\Omega X^{\mathsf T}\ge0$. They act on covariance matrices by $\gamma_B\mapsto
X\gamma_BX^{\mathsf T}+Y$.}
\end{definition}

\begin{proposition}[Properties of channel matrices]\label{prop:ch-matrices}
\reve{\begin{enumerate}[label=(\alph*)]
\item Channel matrices map admissible covariance matrices to admissible ones.
\item Let $L$, $\nu$ and $\gamma$ also denote the real matrices of a linear map and of two bilinear forms
on $\R^{2n}$, and put $X=\Omega^{\mathsf T}L^{\mathsf T}\Omega$ and $Y=\Omega^{\mathsf T}\nu\Omega$. Then
\[\Omega^{\mathsf T}\bigl(L^{\mathsf T}\gamma L+\nu\bigr)\Omega=X\bigl(\Omega^{\mathsf T}\gamma\Omega\bigr)X^{\mathsf
T}+Y .\]
\item For symmetric $\nu$, the matrices $(X,Y)$ of (b) are channel matrices if and only if $L$ and
$\nu$ satisfy \eqref{eq:ch-cond}.
\end{enumerate}}
\end{proposition}

\begin{proof}
\reve{\begin{enumerate}[label=(\alph*)]
\item The output $X\gamma_BX^{\mathsf T}+Y$ is symmetric, and
\[X\gamma_BX^{\mathsf T}+Y-\ii\Omega=X(\gamma_B-\ii\Omega)X^{\mathsf T}+\bigl(Y-\ii\Omega+\ii X\Omega
X^{\mathsf T}\bigr)\]
is a sum of positive semidefinite matrices.
\item Insert $\Omega\Omega^{\mathsf T}=1$ between the factors.
\item The matrices of \eqref{eq:ch-cond} are $U^*HU$ with $U$ real and
$H=\nu+\ii(\Omega-L^{\mathsf T}\Omega L)$, and taking for $U$ the standard basis shows that
\eqref{eq:ch-cond} holds if and only if $H\ge0$. Using $\Omega^{\mathsf T}=-\Omega$ and $\Omega\Omega^{\mathsf
T}=1$ one computes
\[Y-\ii\Omega+\ii X\Omega X^{\mathsf T}=\Omega^{\mathsf T}H^{\mathsf T}\Omega ;\]
congruence by
the orthogonal matrix $\Omega$ and transposition preserve positive semidefiniteness of Hermitian
matrices, and $Y$ is symmetric exactly when $\nu$ is.
\end{enumerate}}
\end{proof}

\noindent\motiv{\reve{For the attenuator, $X=\zeta\cdot1$ and $Y=(1-\zeta^2)1$, and the condition reads
$(1-\zeta^2)(1-\ii\Omega)\ge0$. If channel matrices $(X,Y)$ arise from $L$ and $\nu$ as in (b) and
$Y=0$, then $\nu=0$, and by (c) and Proposition~\ref{prop:ch-special}(a) the map $L$ preserves
$\sigma$. Since $\sigma$ is nondegenerate on $\R^{2n}$, such an $L$ is injective, hence bijective;
then $X=L^{-1}$, \reve{the inverse by which covariance matrices transform
(Corollary~\ref{cor:gs-onsite})}, and by
Proposition~\ref{prop:ch-reversible} the Gaussian channel with these data and $\mu=0$ exists and is
$\alpha_L$.}}

\subsection{Dilations}

\noindent\motiv{\reve{The attenuator arises from a beam splitter: the mode is mixed with an auxiliary
mode prepared in the vacuum, and the auxiliary mode is then discarded. In general, let a second
real vector space $V'$ with an alternating form $\sigma'$ describe an ancilla, and let $S$ be a
linear map of $V\times V'$ preserving the form
\[(\sigma\oplus\sigma')\bigl((u,u'),(w,w')\bigr)=\sigma(u,w)+\sigma'(u',w');\]
for bilinear forms $\gamma$ on $V$ and $\gamma'$ on $V'$, we define the form $\gamma\oplus\gamma'$ in the same way.
When $S$ is bijective, the reversible dynamics $\alpha_S$ of system and ancilla, followed by
discarding an ancilla prepared in a Gaussian state with covariance $\gamma'$, should act on
covariances through the two components of the image of a system vector,
$S(u,0)=(S_{11}u,S_{21}u)$. The following proposition establishes the corresponding statements for
the data; these notes do not carry out the construction at the level of the algebra.}}

\begin{proposition}[Dilation data]\label{prop:ch-dilation}
\reve{Let $S$ preserve $\sigma\oplus\sigma'$, and let $\gamma'$ be a symmetric bilinear form on $V'$ with
$\gamma'+\ii\sigma'\ge0$. Then:
\begin{enumerate}[label=(\alph*)]
\item for all $u,w\in V$,
\[\sigma(S_{11}u,S_{11}w)+\sigma'(S_{21}u,S_{21}w)=\sigma(u,w);\]
\item the data $L=S_{11}$ and $\nu=\gamma'(S_{21}\cdot,S_{21}\cdot)$ satisfy \eqref{eq:ch-cond};
\item for every form $\gamma$ on $V$,
\[(\gamma\oplus\gamma')\bigl(S(u,0),S(w,0)\bigr)=\gamma(S_{11}u,S_{11}w)+\gamma'(S_{21}u,S_{21}w).\]
\end{enumerate}}
\end{proposition}

\begin{proof}
\reve{\begin{enumerate}[label=(\alph*)]
\item This is the preservation of $\sigma\oplus\sigma'$ at $(u,0)$ and $(w,0)$.
\item By (a),
\[\nu+\ii\bigl(\sigma-\sigma(S_{11}\cdot,S_{11}\cdot)\bigr)=(\gamma'+\ii\sigma')(S_{21}\cdot,S_{21}\cdot),\]
so the matrices of \eqref{eq:ch-cond} for a family $(u_j)$ are those of $\gamma'+\ii\sigma'$ for $(S_{21}u_j)$.
\item This is the definition of $\gamma\oplus\gamma'$ at $S(u,0)=(S_{11}u,S_{21}u)$.
\end{enumerate}}
\end{proof}

\noindent\motiv{\reve{Part (c) computes the covariance that the system would have after the dilation,
$\gamma(S_{11}\cdot,S_{11}\cdot)+\gamma'(S_{21}\cdot,S_{21}\cdot)$. By Proposition~\ref{prop:ch-states} this is
the output covariance of the Gaussian channel with $L=S_{11}$ and $\nu=\gamma'(S_{21}\cdot,S_{21}\cdot)$,
whose existence is not proved here, and part (b) says that these data always satisfy the condition
\eqref{eq:ch-cond}. An ancilla state with displacement $d'$ would contribute the displacement
$\mu=d'\circ S_{21}$.}}

\begin{example}[The beam splitter]\label{ex:ch-beamsplitter}
\reve{For one system and one ancilla mode, let $S$ act on positions and on momenta by the same
rotation
\[\begin{pmatrix}\zeta&\sqrt{1-\zeta^2}\\-\sqrt{1-\zeta^2}&\zeta\end{pmatrix},\]
mixing system and ancilla. Since the same orthogonal matrix acts on positions and on momenta, $S$
preserves $\sigma\oplus\sigma'$, and $S_{11}=\zeta\cdot1$,
$S_{21}=-\sqrt{1-\zeta^2}\cdot1$. With the ancilla in the vacuum, $\gamma'(u',w')=u'^{\mathsf T}w'$, the
noise is $(1-\zeta^2)u^{\mathsf T}w$, which is the noise of the attenuator of
Example~\ref{ex:ch-attenuator}. At the level of the algebra, $\alpha_S$ followed by
the vacuum state on the ancilla is expected to give a completely positive map with these data; this
step is not carried out in these notes.}
\end{example}

\subsection{Quantum cellular automata and locality}

\noindent\motiv{\reve{We return to the lattice. Let $G$ be an abelian group and $\cA$ a net on $G$:
each region $O\subseteq G$ is assigned a closed $*$-subalgebra $\cA(O)$ of a C*-algebra $A$,
monotonically in $O$, and a translation action $\alpha$ of $G$ by $*$-automorphisms of $A$ carries
$\cA(O)$ to $\cA(O+a)$ (Chapter~\ref{sec:ccrautomata}). A reversible automaton is a $*$-automorphism commuting
with the translations, together with a finite neighbourhood by which it moves observables at most.
Dropping invertibility, we arrive at the general definition.}}

\begin{definition}[Quantum cellular automaton]\label{def:ch-qca}
\reve{A \emph{quantum cellular automaton} on $\cA$ is a unital completely positive map $T\colon A\to A$
together with a finite set $N\subseteq G$, its \emph{neighbourhood}, such that $T\alpha_a=\alpha_aT$
for all $a\in G$ and $T(\cA(O))\subseteq\cA(O+N)$ for every region $O$.}
\end{definition}

\noindent\motiv{\reve{The last condition bounds how far a single step can spread an observable. A
second notion of locality asks instead that observables in sufficiently separated regions be
processed independently. We state both for an arbitrary linear map.}}

\begin{definition}[Notions of locality]\label{def:ch-locality}
\reve{A linear map $T$ of $A$
\begin{enumerate}[label=(\Roman*)]
\item has \emph{finite propagation with neighbourhood $N$} if $T(\cA(O))\subseteq\cA(O+N)$ for all
regions $O$;
\item \emph{factorizes with respect to $N'\subseteq G$}, where $0\in N'=-N'$, if
$T(b_1b_2)=T(b_1)T(b_2)$ whenever $b_1\in\cA(O_1)$, $b_2\in\cA(O_2)$ and $(O_1+N')\cap O_2=\emptyset$
(so that, since $0\in N'$, the regions are disjoint, and, since $N'=-N'$, the \reve{separation hypothesis}
is symmetric in $O_1$ and $O_2$).
\end{enumerate}}
\end{definition}

\noindent\motiv{\reve{Finite propagation with the neighbourhood $N$ is part of the definition of an
automaton $(T,N)$. Factorization says that the image of a product of observables in sufficiently
separated regions is the product of their images. For the CCR net of Theorem~\ref{thm:ccr-net} the local algebras of disjoint regions commute, and the
product $b_1b_2$ takes the place of the observable $b_1\otimes b_2$ of the composite region. A $*$-automorphism factorizes with respect to every
$N'$, since it is multiplicative; for a channel the product formula of
Proposition~\ref{prop:ch-product} shows what factorization demands of the dissipation form.}}

\begin{proposition}[Reversible automata]\label{prop:ch-revqca}
\reve{If $(T,N)$ is a reversible quantum cellular automaton, then $(T,N)$ is a quantum cellular
automaton.}
\end{proposition}

\begin{proof}
\reve{A $*$-automorphism is unital and completely positive, and the remaining conditions are the
same in both definitions.}
\end{proof}

\subsection{Notes and references}\label{sec:ch-notes}

\begin{enumerate}[label=(\arabic*)]
\item \reve{\emph{Definition.} Definition~\ref{def:ch-channel} is (4.32) of~\cite[p.~92]{krueger},
with the displacement of the standard definition of a quasi-free completely positive map added; the
thesis's linear map and noise are $\Gamma$ and $g$. Condition \eqref{eq:ch-cond} is (4.34). That it
characterizes complete positivity is Theorem~2.6 of~\cite{krueger}, taken from Demoen, Vanheuverzwijn
and Verbeure~\cite{dvv} \vetsource{(scope of the theorem there to be checked)}; these notes assume it
through the definition. The displacement appears already in (2.38) of~\cite{krueger}. \revr{For finitely many modes, the channel matrices of
Definition~\ref{def:ch-matrices} and their condition are equivalent, by conjugation with $\Omega$,
to the phase-space description of Holevo and Werner~\cite{holevo-werner}. In their normalization,
with $\hbar=1$ and the coordinates reordered, the form $\Delta(z,z')$ of (3.12)--(3.13) there is
$-\sigma$, the correlation matrix $\alpha$ of (3.17) is $\gamma/2$, and a Gaussian channel, (4.35)
with (4.45), has $K=L^{\mathsf T}$ and $f(z)=e^{-z^{\mathsf T}\beta z/2}$ with $\beta=\nu/2$, so
that $X=\Omega^{\mathsf T}K\Omega$ and $Y=2\,\Omega^{\mathsf T}\beta\Omega$. For such $f$ the matrices
of their twisted positive definiteness condition (p.~9) are, up to congruence by a positive
diagonal matrix, $[e^{\Delta_T(z_r,z_s)/2}]$; these are positive semidefinite for all families
$(z_r)$ exactly when \eqref{eq:ch-cond} holds, by the Schur product theorem and, conversely, by
adjoining the point $0$ and replacing $z_r$ by $tz_r$ with $t\to0$.} The term ``dissipation form'' for $\Delta_T$ is the thesis's
\reve{\cite[p.~93]{krueger}}.}
\item \reve{\emph{The product formula and states.} Proposition~\ref{prop:ch-product} is (4.35)
of~\cite[p.~92]{krueger}. The thesis prints the output characteristic function (4.33) with $+g(\xi,\xi)/4$;
Proposition~\ref{prop:ch-states} shows that the sign is $-$, in agreement with the covariance update
after (2.41) of~\cite{krueger}. Granted that the attenuator and the vacuum exist, which is not proved
in these notes, the sign $+$ cannot be right: the attenuator would then carry the vacuum to a
functional with $|\chi(u)|=e^{(1-2\zeta^2)u^{\mathsf T}u/4}$, which exceeds $1$ for $u\neq0$ and
$\zeta^2<\frac12$, whereas the characteristic function of a state is bounded by $1$
(Proposition~\ref{prop:gs-twisted}(b)). The two special solutions of (4.34) noted on p.~92 of the thesis are
Proposition~\ref{prop:ch-special}(b) and (c).}
\item \reve{\emph{Dilations.} Proposition~\ref{prop:ch-dilation}(c) is the reduction (4.40)
of~\cite[p.~98]{krueger} for an arbitrary ancilla, and part (b) goes beyond the thesis, which checks
only that the resulting noise is real and symmetric. For finitely many modes, every Gaussian channel
has a Gaussian dilation with a pure Gaussian environment~\cite[Sec.~IV, (30)--(47)]{cegh}.}
\item \reve{\emph{Automata and locality.} Definition~\ref{def:ch-qca} \reve{corresponds to}
Definition~4.1 of~\cite[p.~73]{krueger}\reve{, with the quasi-local algebra there replaced by a net,
as for Definition~\ref{def:rqca}}; the thesis also asks $0\in N$, but by monotonicity of the net, enlarging a
neighbourhood is always allowed, so the condition is not needed here. Notions (I) and (II) of
Definition~\ref{def:ch-locality} are the
first two of the four notions of locality of~\cite[p.~93]{krueger}; the thesis states (II) for
$A_1\otimes A_2$ on disjoint finite regions, with $M=-M$ (and $0\in M$ without loss of generality),
which corresponds to the product $b_1b_2$ here. For a bounded map of the CCR net, asking the condition
for finite regions only is equivalent: every element of the $*$-algebra generated by the Weyl
operators of a region lies in the local algebra of a finite subregion, these elements are dense in
the local algebra, and the product is continuous. In finite dimensions, notion (I) is causality preservation (Definition~1
of~\cite{piroli-cirac}) and notion (II) the factorization condition (6) of~\cite{piroli-cirac}, which
defines their class of factorizing channels (Definition~5 there).}
\end{enumerate}
